\documentclass{article}
\usepackage{default_notes}

\begin{document}

\noindent \colorbox{Apricot}{Conseils pour la dissert}\\

\noindent Mettre en avant les notions qui ne sont pas le thème de 
l'année\\ 
trouver un problème avec la citation \\ 
Mise en page \\
Reprendre les mots exacts du sujet tout au long de la copie
quitte à les redéfinir\\ 
On choisit les citations et on les contextualise \\ 
La truanderie peut fonctionner \\ 
Tout concept doit être définit \\ 

\noindent \Large \colorbox{Aquamarine}{Dissert}\\ 

\noindent \large

L’art n’est pas à mes yeux une réjouissance solitaire. 
Il est un moyen d’émouvoir le plus grand nombre d’hommes en leur offrant 
une image privilégiée des souffrances et des joies communes. Il oblige donc
l’artiste à ne pas se séparer ; il le soumet à la vérité la plus humble et 
la plus universelle. Et celui qui, souvent, a choisi son destin d’artiste
parce qu’il se sentait différent apprend bien vite qu’il ne nourrira son 
art, et sa différence, qu’en avouant sa ressemblance avec tous. 
L’artiste se forge dans cet aller-retour perpétuel de lui aux autres.\\ 
Camus, discours de remise du prix Nobel\\

\noindent
L'art n'est pas une activité individuelle réalisée pour elle-même,
la fonction de l'art est de révèler une expérience 
universelle des hommes. Du coup l'artiste est obligé d'être 
intégré dans la société. Celui qui devient artiste car il se 
sent différent, unique, doit accepter le fait qu'il fait partie d'un 
grand ensemble, qu'il est comme les autres, s'il veut enrichir son art. 
Et en maintenant ce lien avec les autres que l'artiste se construit. \\

Or, rechercher à émouvoir le plus grand nombre d'hommes peut dérailler
le processus créatif et n'encourage pas l'innovation et un 
renouvellement dans l'art. De plus, c'est un reconnaissant ses 
facultés particulières avant tout, que l'artiste se distingue des
autres. Ainsi, en survolant l'individualité de l'artiste 
en faveur de son rôle dans la société, n'ignorons donc pas la nature 
de l'art en tant que médium d'expression personnelle ? 

Alors, on peut interroger le lien qui s'installe entre l'artiste 
et les autres. 
Pour répondre à cette question, à travers l'étude de Ion et de la 
République de Platon, de l'Oeuvre de Zola, et de Un lieu à soi de 
Virginia Woolf, nous montrerons d'abord que l'artiste fait partie 
d'une de la société. Nous verrons ensuite que l'individualité de 
l'artiste est beaucoup plus importante que ses points communs avec 
les autres. Alors, pourrait-on y penser plutôt comment l'expression 
de l'individualité transmet l'universalité  

\noindent 

I \\ 
C'est en admettant sa ressemblance aux autres que l'artiste produit une 
oeuvre qui emeut tout le monde. \\

I-I L'artiste produit ses plus belles oeuvres en dépeignant les 
souffrances et joies communes \\
\(\rightarrow\) Ion\\

I-II L'artiste ne peut pas se séparer complètement des autres, il reste 
un individu qui parcourt la société \\ 

II \\ 
Mais 

II-I L'artiste a besoin de solitude pour créer \\
\(\rightarrow\) Un lieu à soi de Virginia Woolf \\
\(\rightarrow\) L'atelier de Claude\\

II-I Si l'artiste cherche à émouvoir le plus de monde, ne chercherait-il 
pas plus à plaire qu'à vraiment créer une oeuvre d'art ?\\ 
\(\rightarrow\) Flagerolles \\
\(\rightarrow\) C'est là toute la critique de l'art pour Platon  \\


III \\ 
N'est-ce pas en assumant sa différence, son individualité, plutôt que sa 
ressemblance, que l'artiste produit une oeuvre qui émeut 
tout le monde ? \\ 

III-I La solitude offre à l'artiste a une perspective privilégiée 
du monde qui peut permettre de révèler une experience commune dont 
on en a pas conscience\\

III-II C'est en exprimant son individualité à travers son oeuvre 
que l'artiste peut tendre vers l'expression d'une expérience 
universelle \\ 







\end{document}

